\documentclass[10pt, a4paper, twocolumn]{article}

% ── packages ────────────────────────────────────────────────────────────────
\usepackage[a4paper,
            top=22mm, bottom=24mm,
            left=15mm, right=15mm,
            columnsep=7mm]{geometry}
\usepackage[T1]{fontenc}
\usepackage[utf8]{inputenc}
\usepackage[letterspace=150]{microtype}
\usepackage{mathptmx}
\usepackage{parskip}
\usepackage{titlesec}
\usepackage{fancyhdr}
\usepackage{enumitem}
\usepackage{xcolor}
\usepackage{booktabs}
\usepackage{array}
\usepackage{hyperref}
\usepackage{setspace}
\usepackage{amssymb}
\usepackage{authblk}
\usepackage{abstract}
\usepackage{caption}
\usepackage{tikz}
\usetikzlibrary{arrows.meta, positioning, calc}
\usepackage{fontawesome5}

% ── colours ─────────────────────────────────────────────────────────────────
\definecolor{darktext}{HTML}{111111}
\definecolor{mutedgrey}{HTML}{555555}
\definecolor{rulegrey}{HTML}{222222}
% diagram palette
\definecolor{accentblue}{HTML}{1A6FC4}
\definecolor{icongreen}{HTML}{1AAB6D}
\definecolor{iconorange}{HTML}{D4621A}
\definecolor{iconpurple}{HTML}{7B4FBF}
\definecolor{iconred}{HTML}{D44E1A}
\definecolor{labelcol}{HTML}{1A1A2E}
\definecolor{desccol}{HTML}{444455}
\definecolor{arrowcol}{HTML}{BBBBCC}

% ── hyperref ────────────────────────────────────────────────────────────────
\hypersetup{
  colorlinks=false,
  hidelinks,
  pdftitle={Diotima: Compliance-Grade AI Infrastructure for Formative Assessment in Education},
  pdfauthor={Noval Consulting}
}

% ── typography ───────────────────────────────────────────────────────────────
\setlength{\parskip}{4pt}
\setlength{\parindent}{1em}
\frenchspacing

% ── section headings — centred small-caps, like Regional Studies ─────────────
\titleformat{\section}
  {\centering\normalsize\bfseries\scshape}
  {\thesection.}
  {0.5em}
  {}
  []

\titleformat{\subsection}
  {\normalsize\itshape}
  {}
  {0pt}
  {}

\titlespacing{\section}{0pt}{14pt}{6pt}
\titlespacing{\subsection}{0pt}{10pt}{3pt}

% ── abstract styling ─────────────────────────────────────────────────────────
\setlength{\absleftindent}{0pt}
\setlength{\absrightindent}{0pt}
\renewcommand{\abstractnamefont}{\normalsize\bfseries\scshape}
\renewcommand{\abstracttextfont}{\small}

% ── header / footer ──────────────────────────────────────────────────────────
\pagestyle{fancy}
\fancyhf{}
\renewcommand{\headrulewidth}{0pt}
\fancyhead[L]{\small\textit{AI \& Education Policy}}
\fancyhead[R]{\small\textit{Noval Consulting Working Paper}}
\fancyfoot[C]{\small\thepage}
\setlength{\headheight}{14pt}

% ── bullet list style ─────────────────────────────────────────────────────────
\setlist[itemize,1]{
  label={--},
  leftmargin=1.2em,
  itemsep=1pt,
  topsep=3pt,
  parsep=0pt
}

% ── title block — one column, then two-column body ───────────────────────────
\makeatletter
\renewcommand{\maketitle}{%
  \twocolumn[%
    \begin{@twocolumnfalse}
      % journal slug
      {\small\textit{AI \& Education Policy},
       Working Paper Series,\ pp.\ 1--7.}
      \vspace{6pt}
      \hrule height 0.8pt
      \vspace{2pt}
      \hrule height 0.2pt
      \vspace{14pt}

      % title
      {\centering
        {\fontsize{22}{27}\selectfont
         \textbf{Diotima: Compliance-Grade AI Infrastructure}\\[4pt]
         \textbf{for Formative Assessment in Education}}\\[10pt]

        % authors
        {\normalsize\scshape Noval Consulting}\\[4pt]

        % affiliation
        {\small\itshape
          AI Consulting and EU AI Act Advisory Practice,
          Ireland.\quad
          Contact: \texttt{info@novalconsultancy.com}}\\[10pt]
      }

      \hrule height 0.4pt
      \vspace{10pt}

      % abstract
      \begin{abstract}
        \noindent
        The EU AI Act's classification of student assessment systems as high-risk under Annex~III
        creates both a compliance obligation and a need to answer a strategic design question:
        what does an educational AI platform look like when it is built for regulatory alignment
        from first principles, rather than retrofitted to it? This paper describes Diotima, a
        teacher-centred formative assessment platform that treats Annex~III obligations not as
        constraints but as architectural foundations. By grounding all AI-generated content in
        approved curriculum materials, scoping high-risk components narrowly, enforcing human
        oversight as a structural requirement of the system rather than a policy statement, and embedding post-market
        monitoring as an operational discipline, Diotima demonstrates that high-risk AI in
        education can be simultaneously pedagogically sound, technically robust, and legally
        compliant. The paper argues that compliance-by-design is not a cost of participation in
        regulated education markets---it is the only credible path to adoption at scale.
      \end{abstract}

      \vspace{4pt}
      \noindent\textbf{Keywords}\quad
        EU AI Act\quad Annex~III\quad Formative assessment\quad
        Educational technology\quad Human oversight\quad Compliance-by-design

      \vspace{12pt}
      \hrule height 0.8pt
      \vspace{2pt}
      \hrule height 0.2pt
      \vspace{16pt}
    \end{@twocolumnfalse}
  ]%
}
\makeatother

% ─────────────────────────────────────────────────────────────────────────────

\begin{document}

\maketitle
\thispagestyle{fancy}

% ── INTRODUCTION ─────────────────────────────────────────────────────────────
\section{Introduction}

In proposing a new framework for AI governance in education, the EU AI Act draws
a conclusion that practitioners have long recognised implicitly: assessment is not a
neutral activity. \textsc{European Parliament} (2024, Art.~6) classifies AI systems
used to assess students in formal educational settings as high-risk under Annex~III,
Section~3(a). This classification reflects the reality that assessment shapes student
progression and opportunity, that biased or incorrect automated feedback can cause
material harm to learners, and that ungoverned automated evaluation can progressively
displace the professional judgment on which good teaching depends.

Yet the full implications of this regulatory position have not yet been fully absorbed
by the market. Evidence suggests that most educational AI products treat compliance as a
downstream concern---a documentation task to be completed once the product is built,
rather than a constraint that should shape every architectural decision from the outset.
A governance assessment of AI-based assessment systems across Italian universities found
that only 13\% of institutions had adopted AI policies by September~2025, and that
existing documentation prioritised aspirational statements over operational mechanisms
(\textsc{Ferrario et al.}, 2025). At the international level, an OECD review of 18
countries found that none had issued specific regulation on the use of generative AI in
education as of early 2024 (\textsc{OECD}, 2023).

The purpose of this paper is to describe a different approach. Diotima, a
teacher-centred formative assessment platform, serves as a case study for a set of
design principles that are more broadly applicable to any high-risk educational AI
system. Designed from first principles as a compliance-grade system---one that treats
Annex~III obligations not as a constraint on product development but as the
specification for it---Diotima demonstrates that the architecture required for
regulatory alignment is also the architecture required for pedagogically serious,
professionally empowering educational AI. The paper describes Diotima's approach to
human oversight and explainability, its model governance framework, and its data
governance posture, and examines the implications of compliance-by-design for adoption
in regulated education markets.

The paper proceeds as follows. Section~2 establishes the regulatory context and
identifies three foundational challenges that any serious educational AI platform must
address. Section~3 describes Diotima's core design proposition and how it responds to
each of those challenges. Section~4 provides a system overview, illustrated by the lifecycle
diagram in Figure~\ref{fig:lifecycle} and the governance mapping in Table~\ref{tab:mapping}. Section~5 examines the six principal
governance dimensions in turn, presenting each through three lenses: the regulatory
obligation that motivates it, a general design principle applicable to any high-risk
educational AI system, and its specific implementation in Diotima. Section~6 examines
the implications of compliance-by-design for market adoption. Section~7 concludes with
the proposition that the regulatory moment in educational technology is also a design
opportunity---one that rewards systems built for trust rather than retrofitted to it.

% ── REGULATORY CONTEXT ───────────────────────────────────────────────────────
\section{AI in Education Is High-Risk by Default}

The EU AI Act's Annex~III classification of educational assessment AI as high-risk
reflects a straightforward chain of reasoning. Assessment shapes student progression
and opportunity. Biased or incorrect AI outputs can materially harm learners. And if
automated evaluation is allowed to operate without robust human governance, it
progressively displaces the professional judgment that should sit at the centre of any
assessment decision.

The high-risk designation applies even where AI is positioned as formative rather than
summative. Any system that analyses student responses, structures performance
interpretation, and generates feedback can meaningfully influence outcomes---regardless
of how it is labelled. The classification covers not just systems that assign official
grades, but any system whose outputs interact with the assessment process in a way that
could affect how a learner's performance is perceived or recorded (\textsc{European
Parliament}, 2024, Annex~III, s.~3(a)).

Diotima therefore adopts a conservative, regulator-aligned position from the outset: it
is treated as a high-risk system by design. That choice is not defensive---it is
strategic. Every architectural decision, every governance mechanism, every operational
workflow is oriented by that classification.

Three specific challenges follow from this classification that any serious educational
AI platform must address. The first is \textit{hallucination and curriculum drift}: AI
language models can generate plausible but factually incorrect or pedagogically
misaligned outputs that are not inherently grounded in the curriculum being assessed. In formative
assessment contexts, where AI-generated content directly structures how student
performance is interpreted, this is not merely a quality problem but a safety one with
material consequences for learner outcomes. The second is \textit{automation bias and
professional displacement}: well-documented tendencies in human-computer interaction
lead users to accept or defer to automated outputs with insufficient critical scrutiny.
In educational assessment, this risk takes the specific form of teachers progressively
deferring to AI-generated evaluations rather than exercising independent professional
judgment---an outcome that Annex~III's oversight requirements are expressly designed to
prevent. The third is \textit{the governance architecture deficit}: as documented by
\textsc{Ferrario et al.} (2025) and the \textsc{OECD} (2023), the majority of AI
products currently deployed in educational settings were not designed for regulatory
alignment. Compliance has been treated as a retrospective documentation exercise,
creating a structural gap between the obligations that high-risk systems carry and the
architectures that purport to meet them.

% ── CORE PROPOSITION ─────────────────────────────────────────────────────────
\section{Diotima's Core Design Proposition}

The most important thing to establish is what Diotima is \textit{not}. It is not an
AI grader. It does not autonomously assign marks, determine outcomes, or make binding
decisions about student performance. Diotima is a teacher-centred decision-support
system for formative assessment---its function is to expand the scope and consistency
of formative feedback while preserving the judgment calls that professional expertise
requires, not to replace the professional authority at the heart of good teaching.

Each of the three challenges identified in Section~2 has a direct architectural
response. Hallucination and curriculum drift are addressed through mandatory grounding
of all AI-generated content in approved, licensed source materials---the model cannot
generate content that is not explicitly linked to curriculum topics and approved texts.
Automation bias and professional displacement are addressed by making human oversight a
structural requirement of the system rather than a policy aspiration---no AI output is
ever presented to a student without teacher confirmation, and the system is technically
incapable of routing unreviewed outputs to learners. The governance architecture deficit
is addressed by treating Annex~III obligations as the design specification from first
principles, so that compliance evidence is generated as a byproduct of normal operation
rather than assembled retrospectively.

In practice, Diotima supports teachers to: generate curriculum-aligned questions
grounded in approved source materials; apply consistent, transparent rubrics; provide
formative feedback at a scale that would otherwise be impossible; and support student
reflection and iterative improvement rather than delivering a single terminal judgment.

The constraint that structures everything else: no AI output is final without human
confirmation. Teachers approve all assessment content. Teachers validate all evaluative
outputs. The AI proposes; the teacher decides.

% ── SYSTEM OVERVIEW ──────────────────────────────────────────────────────────
\section{System Overview}

Diotima supports the complete formative assessment lifecycle from initial design through
to student reflection and audit. Every stage is engineered to support explainability,
traceability, and human oversight---not as retrospective additions, but as structural
properties of the workflow itself.

The lifecycle proceeds through seven stages: (1)~curriculum-grounded assessment design,
in which teachers select topics, subtopics, and cognitive levels
and the system generates candidate content from approved
source materials; (2)~teacher pre-approval of every generated question, expected answer,
and rubric before any student sees it; (3)~student response and optional AI pre-check, a student-facing feature
that provides provisional analysis to support the learner before final submission; (4)~conditional rubric
visibility, a separate stage in which students receive performance-band information
after final submission, designed to support self-assessment without enabling gaming; (5)~teacher review,
override, and finalisation, during which all AI outputs remain provisional---teachers
may accept, edit, augment, or entirely replace any output before it becomes
authoritative; (6)~student feedback and reflection, in which final
teacher-approved feedback is delivered and students are prompted to engage with their
learning; and (7)~comprehensive logging and auditability, in which every action,
inference, and decision is recorded to support traceability, incident response, and
post-market monitoring.

\begin{figure*}[t]
\centering
\resizebox{0.88\textwidth}{!}{%
\begin{tikzpicture}[
  stagenode/.style={inner sep=0pt, outer sep=0pt},
  connector/.style={->, >=stealth, line width=2.5pt, color=arrowcol,
                    shorten >=2pt, shorten <=2pt},
]
\def\colsep{4.8cm}
%% Row 1
\node[stagenode] (s1) at (0,0) {
  \begin{minipage}{3.4cm}\centering
    {\color{accentblue}\fontsize{28}{28}\selectfont\faBook}\\[5pt]
    {\color{labelcol}\fontsize{9}{9}\selectfont\textbf{\textls[120]{DESIGN}}}\\[4pt]
    {\color{desccol}\fontsize{8}{10}\selectfont
      Curriculum-grounded\\assessment design}
  \end{minipage}};
\node[stagenode] (s2) at (\colsep, 0) {
  \begin{minipage}{3.4cm}\centering
    {\color{icongreen}\fontsize{28}{28}\selectfont\faClipboardCheck}\\[5pt]
    {\color{labelcol}\fontsize{9}{9}\selectfont\textbf{\textls[120]{APPROVE}}}\\[4pt]
    {\color{desccol}\fontsize{8}{10}\selectfont
      Teacher pre-approval of\\all generated content}
  \end{minipage}};
\node[stagenode] (s3) at (2*\colsep, 0) {
  \begin{minipage}{3.4cm}\centering
    {\color{iconorange}\fontsize{28}{28}\selectfont\faPen}\\[5pt]
    {\color{labelcol}\fontsize{9}{9}\selectfont\textbf{\textls[120]{RESPOND}}}\\[4pt]
    {\color{desccol}\fontsize{8}{10}\selectfont
      Student response and\\optional AI pre-check}
  \end{minipage}};
\node[stagenode] (s4) at (3*\colsep, 0) {
  \begin{minipage}{3.4cm}\centering
    {\color{iconpurple}\fontsize{28}{28}\selectfont\faEye}\\[5pt]
    {\color{labelcol}\fontsize{9}{9}\selectfont\textbf{\textls[120]{REVEAL}}}\\[4pt]
    {\color{desccol}\fontsize{8}{10}\selectfont
      Conditional rubric\\visibility for students}
  \end{minipage}};
%% Row 2
\def\rowgap{-3.8cm}
\node[stagenode] (s5) at (3*\colsep, \rowgap) {
  \begin{minipage}{3.4cm}\centering
    {\color{accentblue}\fontsize{28}{28}\selectfont\faSearch}\\[5pt]
    {\color{labelcol}\fontsize{9}{9}\selectfont\textbf{\textls[120]{REVIEW}}}\\[4pt]
    {\color{desccol}\fontsize{8}{10}\selectfont
      Teacher review, override,\\and finalisation}
  \end{minipage}};
\node[stagenode] (s6) at (2*\colsep, \rowgap) {
  \begin{minipage}{3.4cm}\centering
    {\color{icongreen}\fontsize{28}{28}\selectfont\faCommentDots}\\[5pt]
    {\color{labelcol}\fontsize{9}{9}\selectfont\textbf{\textls[120]{REFLECT}}}\\[4pt]
    {\color{desccol}\fontsize{8}{10}\selectfont
      Student feedback and\\reflection}
  \end{minipage}};
\node[stagenode] (s7) at (\colsep, \rowgap) {
  \begin{minipage}{3.4cm}\centering
    {\color{iconred}\fontsize{28}{28}\selectfont\faCheckCircle}\\[5pt]
    {\color{labelcol}\fontsize{9}{9}\selectfont\textbf{\textls[120]{AUDIT}}}\\[4pt]
    {\color{desccol}\fontsize{8}{10}\selectfont
      Comprehensive logging\\and auditability}
  \end{minipage}};
%% Connectors
\draw[connector] ($(s1.east)-(0.1cm,0.55cm)$) -- ($(s2.west)+(0.1cm,-0.55cm)$);
\draw[connector] ($(s2.east)-(0.1cm,0.55cm)$) -- ($(s3.west)+(0.1cm,-0.55cm)$);
\draw[connector] ($(s3.east)-(0.1cm,0.55cm)$) -- ($(s4.west)+(0.1cm,-0.55cm)$);
\draw[connector] ($(s4.south)+(0,-0.05cm)$) -- ($(s5.north)+(0,0.05cm)$);
\draw[connector] ($(s5.west)+(0.1cm,-0.55cm)$) -- ($(s6.east)-(0.1cm,0.55cm)$);
\draw[connector] ($(s6.west)+(0.1cm,-0.55cm)$) -- ($(s7.east)-(0.1cm,0.55cm)$);
\end{tikzpicture}%
}
\caption{The formative assessment lifecycle. Seven stages from curriculum-grounded
  assessment design through to comprehensive logging and audit. Stages~1--4 proceed
  left to right; stages~5--7 return right to left. Every AI-generated output must pass
  through a teacher approval or review gate (stages~2 and~5) before reaching students.}
\label{fig:lifecycle}
\end{figure*}

Figure~\ref{fig:lifecycle} illustrates this lifecycle as a continuous flow from
assessment design through to audit. These lifecycle stages do not each engage a single
governance dimension in isolation. Table~\ref{tab:mapping} maps the six principal
governance dimensions to the lifecycle stages at which each is most directly operative.
The following sections examine each dimension in turn.

\begin{table*}[t]
\centering
\small
\caption{Governance dimensions and lifecycle stage coverage.
  CG~=~Curriculum Grounding; HO~=~Human Oversight; EX~=~Explainability;
  MG~=~Model Governance; DG~=~Data Governance; PM~=~Post-Market Monitoring.
  \checkmark indicates primary operative engagement.}
\label{tab:mapping}
\setlength{\tabcolsep}{5pt}
\begin{tabular}{@{}p{5.8cm}cccccc@{}}
\toprule
\textbf{Lifecycle Stage} & \textbf{CG} & \textbf{HO} & \textbf{EX} & \textbf{MG} & \textbf{DG} & \textbf{PM} \\
\midrule
1.\ Assessment design           & \checkmark &            &            & \checkmark & \checkmark &            \\
2.\ Teacher pre-approval        & \checkmark & \checkmark & \checkmark &            &            &            \\
3.\ Student response / AI check &            &            &            & \checkmark & \checkmark &            \\
4.\ Conditional rubric visibility &          &            & \checkmark &            &            &            \\
5.\ Teacher review and finalisation &        & \checkmark & \checkmark &            &            &            \\
6.\ Student feedback and reflection &        &            & \checkmark &            & \checkmark &            \\
7.\ Logging and auditability    &            &            &            & \checkmark & \checkmark & \checkmark \\
\bottomrule
\end{tabular}
\end{table*}

% ── GOVERNANCE DIMENSIONS IN PRACTICE ─────────────────────────────────────────
\section{Governance Dimensions in Practice}

The following sections examine the six principal governance dimensions that structure
Diotima's compliance posture. Each is presented through three lenses: the
\textit{regulatory obligation} that motivates it, a \textit{general design principle}
applicable to any high-risk educational AI system, and the \textit{specific
implementation in Diotima} that demonstrates the principle in practice. While Diotima
serves as the case study throughout, the design principles described here are not
specific to any single platform---they represent a transferable framework for
compliance-by-design in educational AI.

% ── CURRICULUM GROUNDING ─────────────────────────────────────────────────────
\subsection{Curriculum Grounding and Risk Scoping}

\textit{Regulatory obligation.}\quad Article~9 of the EU AI Act requires high-risk AI
systems to operate within the limits of their intended purpose and to implement risk
management measures that address foreseeable misuse and failure modes. For an
educational AI system, the most direct foreseeable failure mode is the generation of
content that drifts from the curriculum---factually incorrect, pedagogically misaligned,
or ungrounded in the materials that the assessment is intended to evaluate.

\textit{General principle.}\quad High-risk educational AI systems should constrain
generative outputs to approved, verified source materials through retrieval-augmented
or otherwise grounded generation approaches. Content should be anchored not only to
curriculum topics but also to appropriate complexity levels for the target learner
population. The high-risk perimeter should be scoped deliberately to the minimum
necessary components, so that governance obligations remain tractable and risk
propagation is contained.

\textit{In Diotima.}\quad Diotima addresses hallucination architecturally rather
than reactively. Every AI-generated question, answer, and rubric is explicitly linked to
curriculum topics and subtopics, specified learning outcomes, cognitive levels aligned
to the target learner population, and approved, licensed source-of-truth materials.

This grounding serves three functions simultaneously. It prevents the model from
generating content that drifts from the curriculum. It ensures that every item can
be traced back to its source, giving teachers the visibility they need to make informed
approval decisions. And it provides a clear basis for explaining why an item exists and
what it is assessing---which is exactly what regulators and institutions require from
a high-risk system.

A second architectural decision is the deliberate scoping of the high-risk AI
perimeter within the system. Diotima as a whole operates under Annex~III; within
it, the high-risk AI inference is concentrated in two components---the Question,
Answer and Rubric Generation Engine, and the AI Inference and Feedback Engine---while
supporting systems (user interfaces, workflow orchestration, data storage, and
integration components) do not themselves execute high-risk AI functionality and
are governed accordingly. Narrowing the internal high-risk footprint in this way
reduces risk propagation, simplifies governance, and makes compliance obligations
clearer and more tractable.

\textit{Lifecycle coverage.}\quad Grounding operates primarily at stages~1 and~2:
content is anchored to curriculum materials at the point of generation and verified
against those materials during teacher pre-approval.

% ── HUMAN OVERSIGHT ──────────────────────────────────────────────────────────
\subsection{Human Oversight}

\textit{Regulatory obligation.}\quad Article~14 of the EU AI Act requires high-risk AI
systems to be designed and developed so that they can be effectively overseen by natural
persons (i.e.\ human professionals) during their use. The obligation extends to enabling oversight persons to
understand the system's capabilities and limitations, to detect and address failures,
and to intervene on, override, or halt the system's operation.

\textit{General principle.}\quad Human oversight should be enforced as a structural
property of the system workflow, not declared as a policy layer applied after the fact.
The system should be technically incapable of routing AI-generated outputs to affected
persons---in educational contexts, students---without prior review and confirmation by
a qualified human professional. This ensures that accountability remains with the human
decision-maker rather than being diffused across an automated pipeline.

\textit{In Diotima.}\quad Teachers must approve every generated assessment item
before it reaches students. AI rubric placements are always provisional and are never
presented as final. Teachers can accept, edit, override, or entirely replace any AI
output at any stage, and only teacher-confirmed feedback is stored as the authoritative
record. Students never receive AI-generated feedback that a teacher has not reviewed.
That is not a policy aspiration---it is a technical constraint embedded in the system.

The consequences of this design are significant. Professional autonomy is preserved.
The structural risk of over-reliance on AI outputs is prevented rather than mitigated
after the fact. And accountability stays where it belongs---with the human professional,
not the algorithm.

\textit{Lifecycle coverage.}\quad Human oversight is operative primarily at stages~2
and~5: the pre-approval checkpoint before content reaches students, and the review and
finalisation stage before feedback is delivered.

% ── EXPLAINABILITY ───────────────────────────────────────────────────────────
\subsection{Role-Differentiated Explainability}

\textit{Regulatory obligation.}\quad Article~13 of the EU AI Act requires high-risk AI
systems to be sufficiently transparent that users can interpret the system's outputs and
use them appropriately. Transparency obligations extend both to the users operating the
system---in this case, teachers---and, where relevant, to the persons affected by its
outputs---in this case, students.

\textit{General principle.}\quad Explainability should be differentiated by stakeholder
role. System operators---typically teachers or institutional administrators---require
full transparency into the reasoning behind AI outputs, including source materials,
evaluation criteria, and the basis for any automated judgment. Affected
persons---typically students---require contextually appropriate transparency that
supports their engagement with the process without undermining the integrity of the
assessment itself. A single undifferentiated transparency model risks either
overwhelming students with technical detail or depriving professionals of the
information they need to exercise meaningful oversight.

\textit{In Diotima.}\quad For teachers, the system provides full visibility into the
rubric and the source text from which assessment content was derived, so that every
output can be understood in terms of the curriculum materials that grounded it. The pattern of rejection reasons and overrides constitutes an ongoing
signal about model performance---a practical feedback loop that informs model
governance decisions (see Section~\ref{sec:postmarket}).

For students, the explainability design is more deliberate. Rubric information is
revealed conditionally and progressively: students see performance bands after a
genuine attempt, not upfront. Full criteria are not disclosed in advance. This supports
transparency about how responses are being interpreted while preserving the pedagogical
integrity of the assessment, and it orients feedback toward improvement and next steps
rather than terminal judgment.

\textit{Lifecycle coverage.}\quad Explainability operates wherever AI outputs meet
a human: teachers at stages~2 and~5, students at stages~4 and~6.

% ── MODEL GOVERNANCE ─────────────────────────────────────────────────────────
\subsection{Model Governance}

\textit{Regulatory obligation.}\quad Article~9 of the EU AI Act requires high-risk AI
providers to implement risk management measures throughout the system lifecycle,
including identification and analysis of known and foreseeable risks, and estimation
and evaluation of risks that may emerge from actual use. For systems using third-party
foundation models, this obligation extends to the selection, evaluation, and ongoing
monitoring of those models.

\textit{General principle.}\quad For systems built on third-party foundation models,
model selection should be treated as a governance decision, not a technical one. A
structured benchmark suite---covering knowledge and reasoning, factual accuracy,
instruction adherence, bias, and safety---provides the evidential basis for initial
selection and ongoing monitoring. Model versions should be logged with every inference
to support audit, incident investigation, and drift detection. No student data should
be used for model training or fine-tuning.

\textit{In Diotima.}\quad Diotima uses third-party pretrained foundation models
rather than training its own. That is a deliberate and defensible choice---but it does
not mean model selection is a technical decision. In Diotima, model choice is treated
as a governance decision.

Models are evaluated and selected against a structured benchmark suite. Knowledge and
reasoning capabilities are assessed using the Massive Multitask Language Understanding
benchmark (MMLU; \textsc{Hendrycks et al.}, 2021)---a 57-subject evaluation spanning
STEM, humanities, and professional domains---and the Graduate-Level Google-Proof Q\&A
benchmark (GPQA; \textsc{Rein et al.}, 2023), which tests expert-level reasoning across
biology, chemistry, and physics. Factual accuracy and hallucination propensity are
evaluated using LongFact (\textsc{Wei et al.}, 2024), which probes the reliability of
long-form factual claims across a wide range of topics. Instruction-following behaviour
is assessed using the Instruction Following Evaluation benchmark (IFEval;
\textsc{Zhou et al.}, 2023), a verifiable evaluation of adherence to structured natural
language instructions. Bias and fairness are assessed using the Bias Benchmark for
Question Answering (BBQ; \textsc{Parrish et al.}, 2022), a hand-built benchmark
targeting question-answering behaviour across nine social dimensions including age,
gender, and race. Toxicity and safety are assessed using RealToxicityPrompts
(\textsc{Gehman et al.}, 2020), which evaluates the tendency of models to produce
harmful content across a corpus of naturally occurring prompts, and SimpleSafetyTests
(\textsc{Vidgen et al.}, 2023), a test suite of prompts spanning five harm areas
designed to identify critical safety risks such as unsafe advice and toxic content
generation.

No student data is used for model training at any stage. Model versions are logged
with every inference, creating the audit trail needed for incident investigation,
drift detection, and post-market monitoring. This approach enables evidence-based
model selection, ongoing comparison as new models emerge, and safe substitution of
model versions without architectural change.

\textit{Lifecycle coverage.}\quad Model governance is operative at stages~1, 3,
and~7: model selection determines what is generated at the design stage; model
inference is used during the student response and AI pre-check stage; and model
version logging at stage~7 supports post-market monitoring.

% ── DATA GOVERNANCE ──────────────────────────────────────────────────────────
\subsection{Data Governance}

\textit{Regulatory obligation.}\quad Article~10 of the EU AI Act requires high-risk AI
systems to be trained, validated, and tested on data that meets quality criteria
including relevance, representativeness, and freedom from errors and bias. More broadly,
data governance for systems processing personal data must be compatible with applicable
data protection law, including the General Data Protection Regulation (GDPR). For
educational AI systems processing student data---a category that warrants particular
care given the vulnerability of the population---data minimisation and purpose
limitation are foundational requirements.

\textit{General principle.}\quad Data minimisation and purpose limitation should be
foundational design constraints, not compliance additions applied after the fact.
Educational AI systems should process only the data strictly necessary for their stated
purpose, with no repurposing of student data for model training, behavioural profiling,
or secondary commercial use. Where third-party educational content is used as source
material, its governance---including storage, processing, and derivative use---should
be treated as a first-class compliance obligation under explicit licensing agreements.

\textit{In Diotima.}\quad The system processes only what is necessary: student responses,
assessment configuration, and the metadata required to deliver formative feedback and
support audit. There is no collection of sensitive personal categories, no behavioural
profiling, and no use of student data to fine-tune or retrain the underlying models.

Retention limits are defined and aligned with institutional Data Protection Impact
Assessment (DPIA) requirements. Educational content---curriculum documents, licensed
textbooks, and approved source materials---is managed under explicit publisher
agreements covering storage, processing, and derivative use within the scope of
Diotima's purpose.

The effect is a data governance posture that protects every stakeholder simultaneously:
students, whose assessment data is not repurposed; teachers, whose professional actions
are logged for audit but not exposed inappropriately; publishers, whose content is
used under clear and documented licence; and institutions, whose DPIA obligations are
supported rather than complicated by the platform.

\textit{Lifecycle coverage.}\quad Data governance is operative at stages~1, 3, 6,
and~7: at the point of accessing source materials, when processing student responses,
when delivering and storing feedback, and at the comprehensive logging stage.

% ── POST-MARKET MONITORING ───────────────────────────────────────────────────
\subsection{Post-Market Monitoring}
\label{sec:postmarket}

\textit{Regulatory obligation.}\quad Article~72 of the EU AI Act requires providers of
high-risk AI systems to establish and document a post-market monitoring system that
actively collects and analyses data on the system's performance throughout its
operational lifetime. Post-market monitoring is a legal obligation, not a discretionary
quality measure---and it is where many AI governance frameworks, in practice, fail to
deliver on their stated commitments.

\textit{General principle.}\quad Post-market monitoring should be designed so that
evidence of safe operation is generated as a byproduct of normal system use, rather
than requiring separate data collection or retrospective assembly. In educational AI
systems, structured feedback from professional users---teacher rejection reasons,
override patterns, and incident reports---provides a particularly rich source of
real-world performance data, capturing information about model behaviour in authentic
pedagogical contexts that synthetic benchmarks alone cannot replicate.

\textit{In Diotima.}\quad Diotima embeds post-market monitoring from day one as an
operational discipline rather than an afterthought.

Models are regularly re-evaluated against the benchmark suite to detect drift and
performance degradation. Anomaly detection flags distribution shifts in AI outputs
before they reach users. Teachers have structured channels to report harmful or
incorrect outputs, with time-bound remediation workflows. The system is designed for
registration in the EU high-risk AI database, with version update processes in place.

The teacher approval and rejection workflow provides a continuous signal about model
performance in real educational contexts. Rejection reasons and override patterns are
reviewed by human governance personnel and manually entered into the risk register and
model governance process---a structured feedback loop grounded in actual pedagogical
use rather than synthetic evaluation. The system is designed so that this evidence of
safe operation is generated as a byproduct of normal workflow rather than requiring
separate data collection.

\textit{Lifecycle coverage.}\quad Post-market monitoring is operative principally at
stage~7, where comprehensive logging creates the data infrastructure for ongoing
monitoring, but it is also informed continuously by teacher actions at stages~2 and~5.

% ── COMPLIANCE AS FOUNDATION ──────────────────────────────────────────────────
\section{Compliance-by-Design and the Path to Adoption}

The EU AI Act's entry into force marks a structural shift in the conditions for market
participation in educational technology. The obligations imposed by Annex~III---relating
to human oversight, transparency, data governance, model evaluation, and post-market
monitoring---are not documentation requirements. They are design requirements. Products
that treat them as the former face a growing and ultimately unbridgeable gap between
declared and demonstrable compliance.

This gap is not hypothetical. Analysis of the educational AI supplier market suggests
that very few products currently deployed in formal educational settings would meet the
substantive requirements of the Act as currently written---and this pattern is not
confined to the EU. The UK Department for Education published its \textit{Generative AI:
Product Safety Standards} (\textsc{Department for Education}, 2025) setting out the
capabilities and features that AI products must meet to be considered safe for use in
educational settings. Industry analysis by the British Educational Suppliers Association
indicates that as of early 2026, no supplier has been confirmed as fully compliant with
these standards, including products that have received government endorsement (BESA,
personal communication, 2026). The structural reason is the same in both markets:
compliance was not designed in.

Diotima's architecture responds to this directly. Because Annex~III obligations were
treated as the design specification from the outset, the compliance posture is not
asserted---it is verifiable. Each governance dimension documented in this paper
corresponds to a concrete architectural or operational feature that can be examined,
audited, and validated by an institution, a regulator, or a conformity assessment body.
This is what distinguishes compliance-by-design from post-hoc compliance: the former
generates evidence as a byproduct of normal operation; the latter requires it to be
produced retrospectively, from a system not designed to produce it.

The institutional implications are direct. Educational institutions considering
AI-assisted formative assessment do not need to choose between regulatory safety and
pedagogical capability. A system designed from first principles for Annex~III compliance
can be simultaneously safe, transparent, and educationally serious. The architecture
described in this paper demonstrates that compliance, far from constraining what the
technology can do, is the foundation on which trustworthy educational AI must be built.

% ── CONCLUSION ───────────────────────────────────────────────────────────────
\section{Conclusion}

The central proposition of this paper is that the EU AI Act's high-risk classification
of educational assessment AI, rather than constraining what can be built, defines what
must be built. A system that takes Annex~III obligations seriously as design
requirements---grounding its outputs in approved curriculum materials, enforcing human
oversight structurally, differentiating explainability by role, governing model
selection and data use to evidential standards, and embedding post-market monitoring as
an operational discipline---is not a constrained system. It is a more trustworthy, more
auditable, and more professionally empowering one.

Diotima represents one implementation of this proposition. The architecture described
in this paper demonstrates that compliance-by-design is not a theoretical position but
a practically achievable design outcome---one that preserves teacher professional
authority, supports pedagogically serious formative assessment, and generates the
evidence of safe operation that institutions and regulators require.

Architectural foundations, however, do not on their own answer every operational
question. Concrete metrics for measuring compliance in practice---hallucination rates,
validation that teacher review effort is substantive rather than perfunctory, prompt
governance and audit trails, and closed-loop mechanisms that feed teacher corrections
back into model updates---are substantive issues in their own right. A companion paper
will address these operational controls and the measurement frameworks required to
evidence compliance at scale.

The regulatory moment in educational technology creates an obligation that is also an
opportunity: to build AI for education that earns institutional trust rather than
requests it, and that places the professional judgment of teachers at the centre of
every consequential decision about learners.

% ── REFERENCES ───────────────────────────────────────────────────────────────
\section*{References}
\addcontentsline{toc}{section}{References}

\begin{small}
\setlength{\parskip}{3pt}
\noindent
\textsc{Bloom B. S., Engelhart M. D., Furst E. J., Hill W. H.} and \textsc{Krathwohl D. R.}
(1956) \textit{Taxonomy of Educational Objectives: The Classification of Educational
Goals. Handbook I: Cognitive Domain}. David McKay, New York.

\medskip\noindent
\textsc{Department for Education} (2025) Generative AI: Product Safety Standards,
GOV.UK. Available at:
\url{https://www.gov.uk/government/publications/generative-ai-product-safety-standards}
(updated 19~January~2026).

\medskip\noindent
\textsc{European Parliament and Council of the European Union} (2024)
Regulation (EU) 2024/1689 of the European Parliament and of the Council of
13~June~2024 laying down harmonised rules on artificial intelligence (Artificial
Intelligence Act), \textit{Official Journal of the European Union}, L~2024/1689.

\medskip\noindent
\textsc{Ferrario A., Tulli S., Bracco F.} and \textsc{Zanatto D.} (2025)
Testing the applicability of a governance checklist for high-risk AI-based learning
outcome assessment in Italian universities under the EU AI Act Annex~III,
\textit{Frontiers in Artificial Intelligence}, 8, 1718613.

\medskip\noindent
\textsc{Gehman S., Gururangan S., Sap M., Choi Y.} and \textsc{Smith N. A.}
(2020) RealToxicityPrompts: Evaluating Neural Toxic Degeneration in Language
Models, \textit{Findings of the Association for Computational Linguistics:
EMNLP 2020}, pp.\ 3356--3369.

\medskip\noindent
\textsc{Hendrycks D., Burns C., Basart S., Zou A., Mazeika M., Song D.}
and \textsc{Steinhardt J.} (2021)
Measuring Massive Multitask Language Understanding,
\textit{International Conference on Learning Representations (ICLR 2021)},
arXiv:2009.03300.

\medskip\noindent
\textsc{Noval Consulting} (2025)
Diotima AI Assessment Platform: Annex~III High-Risk AI Compliance Dossier, v1.1-draft.
Internal technical documentation, Noval Consulting, Dublin.

\medskip\noindent
\textsc{OECD} (2023) Emerging governance of generative AI in education, in
\textit{OECD Digital Education Outlook 2023}, OECD Publishing, Paris.

\medskip\noindent
\textsc{Parrish A., Chen A., Nangia N., Padmakumar V., Phang J., Thompson J.,
Htut P. M.} and \textsc{Bowman S.} (2022) BBQ: A hand-built bias benchmark for
question answering, \textit{Findings of the Association for Computational
Linguistics: ACL 2022}, pp.\ 2982--2996.

\medskip\noindent
\textsc{Rein D., Hou B. L., Stickland A. C., Petty R., Pang R. Y., Dirani N.,
Michael J.} and \textsc{Bowman S. R.} (2023) GPQA: A Graduate-Level
Google-Proof Q\&A Benchmark,
\textit{arXiv preprint} arXiv:2311.12022.

\medskip\noindent
\textsc{Vidgen B., Thrush T., Waseem Z.} and \textsc{Kiela D.} (2023)
SimpleSafetyTests: a Test Suite for Identifying Critical Safety Risks in Large
Language Models, \textit{arXiv preprint} arXiv:2311.08370.

\medskip\noindent
\textsc{Wei J., Yang C., Song Y., Lu P., Hu E., Tran D., Peng D., Liu R.,
Huang D., Du Y.} and \textsc{Le Q. V.} (2024) Long-form factuality in large
language models,
\textit{arXiv preprint} arXiv:2403.18802.

\medskip\noindent
\textsc{Zhou J., Lu P., Mishra S., Brahma S., Basu S., Luan Y., Zhou D.}
and \textsc{Hou L.} (2023) Instruction-Following Evaluation for Large Language
Models,
\textit{arXiv preprint} arXiv:2311.07911.

\end{small}

\end{document}
